\subsection{Sucesión de Fibonacci}

La sucesión de Fibonacci comienza con 0 y 1, y cada número posterior es la suma de los dos anteriores. Formalmente:
\begin{align*}
    F(0) &= 0 \\
    F(1) &= 1 \\
    F(n) &= F(n-1) + F(n-2) \quad \text{para } n > 1
\end{align*}
El problema radica en calcular el enésimo número de Fibonacci, es decir, $F(n)$, para algún $n$ dado.

\subsubsection{Solución recursiva}

El algoritmo recursivo, que es el más intuitivo, es una traducción directa de la definición matemática:
\begin{pseudocode}
def fibonacci(n: int) -> int:
    if n == 0: return 0
    if n == 1: return 1

    return fibonacci(n-1) + fibonacci(n-2)
\end{pseudocode}

Sin embargo, este algoritmo es ineficiente porque presenta subproblemas superpuestos: para calcular $F(n)$ se necesitan calcular $F(n-1)$ y $F(n-2)$; para calcular $F(n-1)$ se necesitan calcular $F(n-2)$ y $F(n-3)$. En tan solo dos pasos, ya se repitió el cálculo de $F(n-2)$ dos veces. A medida que se avanza, la cantidad de cálculos repetidos crece exponencialmente. De hecho, la complejidad temporal de ese algoritmo es $\Theta(\phi^n)$, donde $\phi$ es el número áureo, aproximadamente 1.618.

Sabiendo que el algoritmo presenta subproblemas superpuestos, solo falta determinar si presenta subestructura óptima para concluir que se puede resolver eficientemente con programación dinámica. En este caso, sí la presenta: la solución óptima para $F(n)$ se construye a partir de las soluciones óptimas para dos subproblemas, $F(n-1)$ y $F(n-2)$.

\subsubsection{Solución top-down}

Una primera optimización es utilizar programación dinámica top-down: se introduce una función principal, \inline{fibonacci_top_down}, que crea un arreglo donde se memoizarán las respuestas. Esa función principal llena el arreglo con $-1$ (que indica que esa solución no se ha calculado) y se lo pasa a la función auxiliar, que memoiza el cálculo si aún no está en el arreglo y retorna la solución memoizada.
\begin{pseudocode}
def fibonacci_top_down(n: int) -> int:
    solutions = |\textrm{arreglo de enteros de tamaño }$n+1$|
    |\br{cbi}|solutions[0] = 0
    solutions[1] = 1|\br{cbf}|

    |\br{lli}|for i=2 to n:
        solutions[i] = -1|\br{llf}|

    return fibonacci_top_down_memoized(n, |\bxl{auxf}|solutions)|\bxr{auxf}|

def fibonacci_top_down_memoized(n: int, solutions: arreglo[int]) -> int:
    |\bxl{memo}|if solutions[n] < 0:|\bxr{memo}|
        solutions[n] = (
            fibonacci_top_down_memoized(n - 1, solutions) + 
            fibonacci_top_down_memoized(n - 2, solutions)
        )
    |\bxl{rmemo}|return solutions[n]|\bxr{rmemo}|
\end{pseudocode}
\begin{annotations}
    \codebrace{cbi}{cbf}{Se añaden los casos base directamente al arrreglo}
    \codebrace[width=9cm]{lli}{llf}{Se llena el resto del arreglo con -1, que indica que no se ha calculado ese término de la secuencia. Es inambiguo porque los términos calculados serán positivos}
    \codebox[width=4.5cm]{auxf}{La función auxiliar recursiva recibe siempre la misma estructura y la va llenando}
    \codebox[width=10cm]{memo}{La estructura aún no contiene la solución para este valor, entonces, se memoiza}
    \codebox[width=6.5cm]{rmemo}{Se retorna siempre el valor memoizado (que puede que se haya calculado en ese llamado recursivo o en uno anterior)}
\end{annotations}

Por definición, para calcular el enésimo número de Fibonacci es necesario calcular todos los anteriores, por ende, el algoritmo top-down es menos eficiente que el algoritmo bottom-up: ambos calculan todos los términos anteriores, pero el bottom-up evita el uso de recursión.

\subsubsection{Solución bottom-up}

El método con programación dinámica bottom-up resuelve primero los problemas más pequeños. En este caso, eso traduce naturalmente a calcular primero los menores términos de la sucesión, o lo que es lo mismo, a calcularlos en orden:

\begin{pseudocode}
def fibonacci(n: int) -> int:
    solutions = |\textrm{arreglo de }|int|\textrm{ de tamaño }$n+1$|
    solutions[0] = 0
    solutions[1] = 1

    for i=2 to n:
        solutions[i] = solutions[i-1] + solutions[i-2]

    return solutions[n]
\end{pseudocode}
Como optimización adicional, este problema tiene la particularidad de que no hace falta preservar las soluciones previas, de forma que se puede ahorrar el espacio en memoria que ocupa el arreglo de soluciones:
\begin{pseudocode}
def fibonacci(n: int) -> int:
    if n <= 1: return n

    prev_prev_solution = 0
    prev_solution = 1

    for i=2 to n:
        solution = prev_solution + prev_prev_solution
        prev_prev_solution = prev_solution
        prev_solution = solution

    return solution
\end{pseudocode}
Se obtiene un algoritmo cuya complejidad temporal es $\Theta(n)$ y cuya complejidad espacial es $\Theta(1)$.

\practicar[leetcode]{Fibonacci Number}{https://leetcode.com/problems/fibonacci-number/}
